\scanpage{115}

% manuscript lines / position: 1–3, crossed block
\deleted{Sous réserve que les $\alpha_i:A_i\unsure{\longrightarrow|\mathord{-}} A$ soient bien injectifs, cela montre que chacun des $T_i$ est déterminé par les deux autres, via ``}

% manuscript lines / position: 4, crossed block
\deleted{\[
 \alpha_1T_1(x)=2x+\deleted{\unsure{\alpha_2T_2(x)-\alpha_3T_3}}\inserted{T_{A/k}(x)-\alpha_2T_2(x)-\alpha_3T_3(x)}
\]}

% manuscript lines / position: 5–6, crossed block
\deleted{On s'attend à trouver \deleted{des} formules similaires, multiplicatives}

% manuscript lines / position: 7, crossed block
\deleted{\[
 \boxed{N_{A/k}(x)\,x^2=N_1(x)N_2(x)N_3(x)}
\]}

% manuscript lines / position: 8, crossed block
\deleted{\[
 \text{et}\qquad \deleted{\unsure{N_2}}\,\boxed{n_1(N_1(x))=n_2(N_2(x))=n_3(N_3(x))},
\]}

% manuscript lines / position: 9–11
Mais \deleted{toutes} les \deleted{conditions} \inserted{relations sur la norme} se prouvent aisément, comme \uline{conséquences} de celles sur les traces, i.e. la condition $\sigma_1\sigma_2\sigma_3=1$ sur les $\sigma_i$. En effet, posons

% manuscript lines / position: 12
\[
 x_0=x,\quad x_1=T_1(x),\quad x_2=T_2(x),\quad x_3=T_3(x)
\]

% manuscript lines / position: 13–15
On a alors
\diagram{\(\begin{array}{lll}
 \sigma_1\unsure{:} & x_0\Leftrightarrow x_1 & x_2\leftrightarrow x_3\\
 \sigma_2: & x_0\leftrightarrow x_2 & x_1\unsure{\longleftrightarrow|\longrightarrow}x_3\\
 \sigma_3: & x_0\leftrightarrow x_3, & x_1\unsure{\longleftarrow|\longleftrightarrow}x_2
 \end{array}\)}

% manuscript lines / position: 16–18
\[
 \begin{aligned}
 \text{d'où}\quad &\deleted{N_1(N_1(x))=N_1(x)\overline{N_1(x)}=x\unsure{\sigma_1(x)}(\unsure{x\sigma_1(x)})}\\
 &\deleted{\unsure{N_1(x_1)=x_1\bar x_1}}\\
 n_1(N_1(x))&\overset{!}{=}\underset{\substack{\Vert\\x\deleted{\unsure{T}}\sigma_1(x)}}{\underbrace{N_1(x)}}\,\overline{N_1(x)}
 =(x_0x_1)\underbrace{\sigma_2(x_0x_1)}_{x_2x_3}=x_0x_1x_2x_3
 \end{aligned}
\]

% manuscript lines / position: 19–21
et de $\unsure{=|\text{même}}$ pour $n_2(N_2(x))$, $n_3(N_3(x))$. Donc on a l'égalité des $n_i(N_i(x))$ si $k\unsure{\mathord{-}|\longrightarrow}A$ est \uline{injectif}. Et on aura

% manuscript lines / position: 22–23
\[
 \begin{aligned}
 N_1(x)N_2(x)N_3(x)&=(x_0x_1)(x_0x_2)(x_0x_3)\\
 &=x_0^2\underbrace{(x_0x_1x_2x_3)}_{N_{A/k}(\unsure{x_0|x})}=x^{\unsure{2}}N_{A/k}(x),
 \end{aligned}
\]

% manuscript lines / position: 24
O.K.

% manuscript lines / position: bottom left
% archive sheet number: 114
